% TMLR submission.
%
% To build the submission version, drop the official tmlr.sty (and tmlr.bst)
% into this directory; the preamble below picks it up automatically. Without
% it the file still compiles standalone via the article fallback, which is how
% it is checked in CI.
%
% Every numeric claim resolves through a macro defined in numbers.tex, which is
% generated by scripts/export_paper_numbers.py from the committed JSON
% artifacts. Do not type a number into this file.

% TMLR requires its own stylefile: "Any changes to the stylefile or template
% that alters the formatting, font, or layout of the manuscript may result in
% rejection without review." Fetch it from
% https://github.com/JmlrOrg/tmlr-style-file and place tmlr.sty and tmlr.bst
% here before submitting. Without them this falls back to article so the draft
% still builds in CI; the fallback deliberately reproduces the anonymous title
% block so a local preview matches what reviewers see.
\newif\ifusetmlr
\IfFileExists{tmlr.sty}{\usetmlrtrue}{\usetmlrfalse}

\ifusetmlr
  \documentclass[10pt]{article}
  \usepackage{tmlr}          % submission: anonymous, "Under review as submission to TMLR"
  % \usepackage[accepted]{tmlr}   % camera-ready
  % \usepackage[preprint]{tmlr}   % de-anonymized preprint (arXiv)
\else
  \documentclass[10pt]{article}
  \usepackage[letterpaper,margin=1.1in]{geometry}
  \usepackage{times}
\fi

\usepackage{amsmath}
\usepackage{amssymb}
\usepackage{booktabs}
\usepackage{graphicx}
\usepackage{microtype}
\usepackage{url}
\usepackage[hidelinks]{hyperref}
\usepackage[round]{natbib}

% Generated by scripts/export_paper_numbers.py -- do not edit by hand.
% Every numeric claim in the manuscript resolves through these macros.

\newcommand{\AOneCollMaxName}{GHL}
\newcommand{\AOneCollMaxRate}{0.5876}
\newcommand{\AOneCollMinName}{GECCO}
\newcommand{\AOneCollMinRate}{0.1733}
\newcommand{\AlphaMSL}{1.0000}
\newcommand{\AlphaPSM}{0.4861}
\newcommand{\AlphaSMAP}{1.0000}
\newcommand{\AlphaSMD}{0.6147}
\newcommand{\AlphaSWaT}{1.0000}
\newcommand{\AlphaWADI}{1.0000}
\newcommand{\ClsCCFlips}{3}
\newcommand{\ClsCCPairs}{6}
\newcommand{\ClsCCRate}{0.5000}
\newcommand{\ClsDCFlips}{27}
\newcommand{\ClsDCPairs}{60}
\newcommand{\ClsDCRate}{0.4500}
\newcommand{\ClsDDFlips}{14}
\newcommand{\ClsDDPairs}{60}
\newcommand{\ClsDDRate}{0.2333}
\newcommand{\ColAlphaP}{0.827}
\newcommand{\ColAlphaRho}{0.0563}
\newcommand{\ColDensP}{0.195}
\newcommand{\ColDensRho}{-0.3333}
\newcommand{\ColDurP}{0.138}
\newcommand{\ColDurRho}{0.3750}
\newcommand{\ColLenP}{0.289}
\newcommand{\ColLenRho}{0.2696}
\newcommand{\ColNsegP}{0.098}
\newcommand{\ColNsegRho}{-0.4118}
\newcommand{\CollMaxName}{GHL}
\newcommand{\CollMaxRate}{0.5876}
\newcommand{\CollMinName}{TAO}
\newcommand{\CollMinRate}{0.1311}
\newcommand{\ConstAlpha}{16}
\newcommand{\ConstBigAlpha}{7}
\newcommand{\ConstBigDur}{0}
\newcommand{\ConstBigNseg}{0}
\newcommand{\ConstDur}{5}
\newcommand{\ConstMultiAlpha}{11}
\newcommand{\ConstMultiDur}{0}
\newcommand{\ConstMultiNseg}{0}
\newcommand{\ConstNseg}{5}
\newcommand{\DeepMSLFlips}{0}
\newcommand{\DeepMSLPairs}{10}
\newcommand{\DeepPSMFlips}{6}
\newcommand{\DeepPSMPairs}{10}
\newcommand{\DeepSMAPFlips}{1}
\newcommand{\DeepSMAPPairs}{10}
\newcommand{\DeepSMDFlips}{1}
\newcommand{\DeepSMDPairs}{10}
\newcommand{\DeepSWaTFlips}{5}
\newcommand{\DeepSWaTPairs}{10}
\newcommand{\DeepWADIFlips}{1}
\newcommand{\DeepWADIPairs}{10}
\newcommand{\DeepWithIForestFlips}{16}
\newcommand{\DeepWithIForestPairs}{30}
\newcommand{\DeepWithIForestRate}{0.5333}
\newcommand{\DeepWithLOFFlips}{11}
\newcommand{\DeepWithLOFPairs}{30}
\newcommand{\DeepWithLOFRate}{0.3667}
\newcommand{\IdentAlphaOne}{0.0}
\newcommand{\IdentAlphaZero}{0.0}
\newcommand{\JointCiHiOneZero}{0.0954}
\newcommand{\JointCiHiTwoZero}{0.0417}
\newcommand{\JointCiHiZeroFive}{0.1783}
\newcommand{\JointCiLoOneZero}{0.0373}
\newcommand{\JointCiLoTwoZero}{0.0092}
\newcommand{\JointCiLoZeroFive}{0.0905}
\newcommand{\JointFirstCollFlipsOneZero}{33}
\newcommand{\JointFirstCollFlipsTwoZero}{49}
\newcommand{\JointFirstCollFlipsZeroFive}{25}
\newcommand{\JointFirstCollOneZero}{CATSv2}
\newcommand{\JointFirstCollPairsOneZero}{213}
\newcommand{\JointFirstCollPairsTwoZero}{675}
\newcommand{\JointFirstCollPairsZeroFive}{79}
\newcommand{\JointFirstCollRateOneZero}{0.1549}
\newcommand{\JointFirstCollRateTwoZero}{0.0726}
\newcommand{\JointFirstCollRateZeroFive}{0.3165}
\newcommand{\JointFirstCollTwoZero}{SMAP}
\newcommand{\JointFirstCollZeroFive}{Genesis}
\newcommand{\JointGhlFlipsOneZero}{5}
\newcommand{\JointGhlFlipsTwoZero}{0}
\newcommand{\JointGhlFlipsZeroFive}{26}
\newcommand{\JointGhlPairsOneZero}{156}
\newcommand{\JointGhlPairsTwoZero}{63}
\newcommand{\JointGhlPairsZeroFive}{294}
\newcommand{\JointNCollOneZero}{17}
\newcommand{\JointNCollTwoZero}{16}
\newcommand{\JointNCollZeroFive}{17}
\newcommand{\JointNullOneZero}{0.4988}
\newcommand{\JointNullTwoZero}{0.4991}
\newcommand{\JointNullZeroFive}{0.5014}
\newcommand{\JointSecondCollFlipsOneZero}{167}
\newcommand{\JointSecondCollFlipsTwoZero}{29}
\newcommand{\JointSecondCollFlipsZeroFive}{158}
\newcommand{\JointSecondCollOneZero}{MSL}
\newcommand{\JointSecondCollPairsOneZero}{1{,}203}
\newcommand{\JointSecondCollPairsTwoZero}{492}
\newcommand{\JointSecondCollPairsZeroFive}{548}
\newcommand{\JointSecondCollRateOneZero}{0.1388}
\newcommand{\JointSecondCollRateTwoZero}{0.0589}
\newcommand{\JointSecondCollRateZeroFive}{0.2883}
\newcommand{\JointSecondCollTwoZero}{MSL}
\newcommand{\JointSecondCollZeroFive}{CATSv2}
\newcommand{\JointSeriesOneZero}{151}
\newcommand{\JointSeriesTwoZero}{128}
\newcommand{\JointSeriesZeroFive}{161}
\newcommand{\JointZeroCollOneZero}{2}
\newcommand{\JointZeroCollTwoZero}{9}
\newcommand{\JointZeroCollZeroFive}{0}
\newcommand{\LocoAlphaMax}{0.3741}
\newcommand{\LocoAlphaMin}{0.0898}
\newcommand{\LocoAlphaWorst}{TAO}
\newcommand{\LocoDensMax}{-0.2144}
\newcommand{\LocoDensMin}{-0.4088}
\newcommand{\LocoDensWorst}{GHL}
\newcommand{\LocoDurMax}{0.4814}
\newcommand{\LocoDurMin}{0.2408}
\newcommand{\LocoDurWorst}{GHL}
\newcommand{\LocoLenMax}{0.3452}
\newcommand{\LocoLenMin}{0.0394}
\newcommand{\LocoLenWorst}{GHL}
\newcommand{\LocoNsegMax}{-0.2535}
\newcommand{\LocoNsegMin}{-0.4252}
\newcommand{\LocoNsegWorst}{TAO}
\newcommand{\ModelATFlips}{15}
\newcommand{\ModelATPairs}{36}
\newcommand{\ModelATRate}{0.4167}
\newcommand{\ModelIForestFlips}{19}
\newcommand{\ModelIForestPairs}{36}
\newcommand{\ModelIForestRate}{0.5278}
\newcommand{\ModelLOFFlips}{14}
\newcommand{\ModelLOFPairs}{36}
\newcommand{\ModelLOFRate}{0.3889}
\newcommand{\ModelMOMENTFlips}{8}
\newcommand{\ModelMOMENTPairs}{36}
\newcommand{\ModelMOMENTRate}{0.2222}
\newcommand{\NBigCollections}{8}
\newcommand{\NMultiCollections}{12}
\newcommand{\NrCiHiOneZero}{0.6155}
\newcommand{\NrCiHiZeroFive}{0.5410}
\newcommand{\NrCiHiZeroTwo}{0.4923}
\newcommand{\NrCiHiZeroZero}{0.4921}
\newcommand{\NrCiLoOneZero}{0.1765}
\newcommand{\NrCiLoZeroFive}{0.2000}
\newcommand{\NrCiLoZeroTwo}{0.1765}
\newcommand{\NrCiLoZeroZero}{0.2143}
\newcommand{\NrClsOneZero}{5}
\newcommand{\NrClsZeroFive}{5}
\newcommand{\NrClsZeroTwo}{4}
\newcommand{\NrClsZeroZero}{0}
\newcommand{\NrDeepOneZero}{16}
\newcommand{\NrDeepZeroFive}{9}
\newcommand{\NrDeepZeroTwo}{5}
\newcommand{\NrDeepZeroZero}{0}
\newcommand{\NrFlipsOneZero}{13}
\newcommand{\NrFlipsZeroFive}{20}
\newcommand{\NrFlipsZeroTwo}{24}
\newcommand{\NrFlipsZeroZero}{44}
\newcommand{\NrPairsOneZero}{35}
\newcommand{\NrPairsZeroFive}{53}
\newcommand{\NrPairsZeroTwo}{75}
\newcommand{\NrPairsZeroZero}{126}
\newcommand{\NrRateOneZero}{0.3714}
\newcommand{\NrRateZeroFive}{0.3774}
\newcommand{\NrRateZeroTwo}{0.3200}
\newcommand{\NrRateZeroZero}{0.3492}
\newcommand{\SerAlphaP}{0.0002}
\newcommand{\SerAlphaRho}{0.3241}
\newcommand{\SerDensP}{0.0001}
\newcommand{\SerDensRho}{-0.3601}
\newcommand{\SerDurP}{0.0001}
\newcommand{\SerDurRho}{0.3717}
\newcommand{\SerLenP}{0.0051}
\newcommand{\SerLenRho}{0.2047}
\newcommand{\SerNsegP}{0.0001}
\newcommand{\SerNsegRho}{-0.3775}
\newcommand{\StratAffAMid}{0.2897}
\newcommand{\StratAffAMidFlips}{1{,}280}
\newcommand{\StratAffAMidPairs}{4{,}419}
\newcommand{\StratAffAMidSeries}{15}
\newcommand{\StratAffAOne}{0.3321}
\newcommand{\StratAffAOneFlips}{13{,}349}
\newcommand{\StratAffAOnePairs}{40{,}200}
\newcommand{\StratAffAOneSeries}{153}
\newcommand{\StratAffAZero}{0.1466}
\newcommand{\StratAffAZeroFlips}{522}
\newcommand{\StratAffAZeroPairs}{3{,}561}
\newcommand{\StratAffAZeroSeries}{12}
\newcommand{\StratSaeAMid}{0.1758}
\newcommand{\StratSaeAMidFlips}{791}
\newcommand{\StratSaeAMidPairs}{4{,}500}
\newcommand{\StratSaeAMidSeries}{15}
\newcommand{\StratSaeAOne}{0.3321}
\newcommand{\StratSaeAOneFlips}{13{,}349}
\newcommand{\StratSaeAOnePairs}{40{,}200}
\newcommand{\StratSaeAOneSeries}{153}
\newcommand{\StratSaeAZero}{0.0000}
\newcommand{\StratSaeAZeroFlips}{0}
\newcommand{\StratSaeAZeroPairs}{3{,}594}
\newcommand{\StratSaeAZeroSeries}{12}
\newcommand{\TabAllAffAgreePct}{65.1}
\newcommand{\TabAllAffCiClusterHi}{0.4841}
\newcommand{\TabAllAffCiClusterLo}{0.2143}
\newcommand{\TabAllAffCiPairHi}{0.4286}
\newcommand{\TabAllAffCiPairLo}{0.2698}
\newcommand{\TabAllAffFlips}{44}
\newcommand{\TabAllAffNull}{0.5021}
\newcommand{\TabAllAffNullSd}{0.0638}
\newcommand{\TabAllAffPairs}{126}
\newcommand{\TabAllAffRate}{0.3492}
\newcommand{\TabAllSaeAgreePct}{71.4}
\newcommand{\TabAllSaeCiClusterHi}{0.4206}
\newcommand{\TabAllSaeCiClusterLo}{0.1508}
\newcommand{\TabAllSaeCiPairHi}{0.3651}
\newcommand{\TabAllSaeCiPairLo}{0.2063}
\newcommand{\TabAllSaeFlips}{36}
\newcommand{\TabAllSaeNull}{0.4991}
\newcommand{\TabAllSaeNullSd}{0.0644}
\newcommand{\TabAllSaePairs}{126}
\newcommand{\TabAllSaeRate}{0.2857}
\newcommand{\TabAlphaLtOne}{2}
\newcommand{\TabDatasets}{6}
\newcommand{\TabDeepAffAgreePct}{76.7}
\newcommand{\TabDeepAffCiClusterHi}{0.4167}
\newcommand{\TabDeepAffCiClusterLo}{0.0667}
\newcommand{\TabDeepAffCiPairHi}{0.3500}
\newcommand{\TabDeepAffCiPairLo}{0.1333}
\newcommand{\TabDeepAffFlips}{14}
\newcommand{\TabDeepAffNull}{0.5004}
\newcommand{\TabDeepAffNullSd}{0.0845}
\newcommand{\TabDeepAffPairs}{60}
\newcommand{\TabDeepAffRate}{0.2333}
\newcommand{\TabDeepSaeAgreePct}{86.7}
\newcommand{\TabDeepSaeCiClusterHi}{0.2833}
\newcommand{\TabDeepSaeCiClusterLo}{0.0333}
\newcommand{\TabDeepSaeCiPairHi}{0.2333}
\newcommand{\TabDeepSaeCiPairLo}{0.0500}
\newcommand{\TabDeepSaeFlips}{8}
\newcommand{\TabDeepSaeNull}{0.5018}
\newcommand{\TabDeepSaeNullSd}{0.0834}
\newcommand{\TabDeepSaePairs}{60}
\newcommand{\TabDeepSaeRate}{0.1333}
\newcommand{\TsbadAgreePct}{68.6}
\newcommand{\TsbadCiClusterHi}{0.3715}
\newcommand{\TsbadCiClusterLo}{0.2695}
\newcommand{\TsbadCiSeriesHi}{0.3312}
\newcommand{\TsbadCiSeriesLo}{0.2979}
\newcommand{\TsbadCollections}{17}
\newcommand{\TsbadDropShareZeroFive}{70.2}
\newcommand{\TsbadDropShareZeroOne}{90.2}
\newcommand{\TsbadDropShareZeroTwo}{83.7}
\newcommand{\TsbadDropZeroFive}{0.2556}
\newcommand{\TsbadDropZeroOne}{0.2972}
\newcommand{\TsbadDropZeroTwo}{0.2827}
\newcommand{\TsbadFlipAff}{0.3145}
\newcommand{\TsbadFlipAffFlips}{15{,}151}
\newcommand{\TsbadFlipAffPairs}{48{,}180}
\newcommand{\TsbadFlipSae}{0.2928}
\newcommand{\TsbadGapOneZeroToTwoZero}{0.2603}
\newcommand{\TsbadGapPairsOneZeroToTwoZero}{10{,}026}
\newcommand{\TsbadGapPairsTwoZeroToInf}{16{,}045}
\newcommand{\TsbadGapPairsZeroFiveToOneZero}{7{,}757}
\newcommand{\TsbadGapPairsZeroOneToZeroTwo}{3{,}139}
\newcommand{\TsbadGapPairsZeroTwoToZeroFive}{6{,}513}
\newcommand{\TsbadGapPairsZeroZeroToZeroOne}{4{,}700}
\newcommand{\TsbadGapShareOneZeroToTwoZero}{20.8}
\newcommand{\TsbadGapShareTwoZeroToInf}{33.3}
\newcommand{\TsbadGapShareZeroFiveToOneZero}{16.1}
\newcommand{\TsbadGapShareZeroOneToZeroTwo}{6.5}
\newcommand{\TsbadGapShareZeroTwoToZeroFive}{13.5}
\newcommand{\TsbadGapShareZeroZeroToZeroOne}{9.8}
\newcommand{\TsbadGapTwoZeroToInf}{0.2015}
\newcommand{\TsbadGapZeroFiveToOneZero}{0.3616}
\newcommand{\TsbadGapZeroOneToZeroTwo}{0.4839}
\newcommand{\TsbadGapZeroTwoToZeroFive}{0.4230}
\newcommand{\TsbadGapZeroZeroToZeroOne}{0.4743}
\newcommand{\TsbadJointOneZero}{0.0651}
\newcommand{\TsbadJointPairsOneZero}{12{,}796}
\newcommand{\TsbadJointPairsTwoZero}{6{,}412}
\newcommand{\TsbadJointPairsZeroFive}{20{,}214}
\newcommand{\TsbadJointShareOneZero}{26.6}
\newcommand{\TsbadJointShareTwoZero}{13.3}
\newcommand{\TsbadJointShareZeroFive}{42.0}
\newcommand{\TsbadJointTwoZero}{0.0220}
\newcommand{\TsbadJointZeroFive}{0.1306}
\newcommand{\TsbadModels}{25}
\newcommand{\TsbadNull}{0.5002}
\newcommand{\TsbadNullPerm}{200}
\newcommand{\TsbadNullSd}{0.0056}
\newcommand{\TsbadRows}{4{,}498}
\newcommand{\TsbadSecFourZeroToInf}{0.0848}
\newcommand{\TsbadSecOneZeroToTwoZero}{0.2272}
\newcommand{\TsbadSecPairsFourZeroToInf}{5{,}247}
\newcommand{\TsbadSecPairsOneZeroToTwoZero}{7{,}628}
\newcommand{\TsbadSecPairsTwoZeroToFourZero}{6{,}417}
\newcommand{\TsbadSecPairsZeroFiveToOneZero}{7{,}404}
\newcommand{\TsbadSecPairsZeroTwoToZeroFive}{8{,}601}
\newcommand{\TsbadSecPairsZeroZeroToZeroTwo}{12{,}883}
\newcommand{\TsbadSecTwoZeroToFourZero}{0.1733}
\newcommand{\TsbadSecZeroFiveToOneZero}{0.3052}
\newcommand{\TsbadSecZeroTwoToZeroFive}{0.3675}
\newcommand{\TsbadSecZeroZeroToZeroTwo}{0.4999}
\newcommand{\TsbadSeries}{180}
\newcommand{\TwoColAlphaPOneZero}{0.476}
\newcommand{\TwoColAlphaPTwoZero}{0.899}
\newcommand{\TwoColAlphaPZeroFive}{0.631}
\newcommand{\TwoColAlphaRhoOneZero}{0.1881}
\newcommand{\TwoColAlphaRhoTwoZero}{0.0362}
\newcommand{\TwoColAlphaRhoZeroFive}{0.1279}
\newcommand{\TwoColDurPOneZero}{0.295}
\newcommand{\TwoColDurPTwoZero}{0.976}
\newcommand{\TwoColDurPZeroFive}{0.521}
\newcommand{\TwoColDurRhoOneZero}{0.2685}
\newcommand{\TwoColDurRhoTwoZero}{-0.0097}
\newcommand{\TwoColDurRhoZeroFive}{0.1667}
\newcommand{\TwoColNsegPOneZero}{0.068}
\newcommand{\TwoColNsegPTwoZero}{0.425}
\newcommand{\TwoColNsegPZeroFive}{0.008}
\newcommand{\TwoColNsegRhoOneZero}{-0.4574}
\newcommand{\TwoColNsegRhoTwoZero}{-0.2139}
\newcommand{\TwoColNsegRhoZeroFive}{-0.6225}
\newcommand{\TwoSerAlphaPOneZero}{0.433}
\newcommand{\TwoSerAlphaPTwoZero}{0.859}
\newcommand{\TwoSerAlphaPZeroFive}{0.291}
\newcommand{\TwoSerAlphaRhoOneZero}{0.0642}
\newcommand{\TwoSerAlphaRhoTwoZero}{0.0159}
\newcommand{\TwoSerAlphaRhoZeroFive}{0.0858}
\newcommand{\TwoSerDurPOneZero}{0.423}
\newcommand{\TwoSerDurPTwoZero}{0.447}
\newcommand{\TwoSerDurPZeroFive}{0.946}
\newcommand{\TwoSerDurRhoOneZero}{0.0666}
\newcommand{\TwoSerDurRhoTwoZero}{0.0687}
\newcommand{\TwoSerDurRhoZeroFive}{-0.0056}
\newcommand{\TwoSerNsegPOneZero}{0.000}
\newcommand{\TwoSerNsegPTwoZero}{0.220}
\newcommand{\TwoSerNsegPZeroFive}{0.000}
\newcommand{\TwoSerNsegRhoOneZero}{-0.3193}
\newcommand{\TwoSerNsegRhoTwoZero}{-0.1098}
\newcommand{\TwoSerNsegRhoZeroFive}{-0.3294}


\newcommand{\afff}{\mbox{Aff-F1}}
\newcommand{\auc}{\mbox{AUC-ROC}}
\newcommand{\sae}{\mbox{SAEScore}}

\title{How Much Do Time-Series Anomaly Detection Metrics Actually Disagree?\\
Null Models and Cluster-Aware Inference for Rank-Flip Statistics}

% TMLR is double blind and rejects non-anonymous submissions without review.
% tmlr.sty suppresses this block automatically unless [accepted] or [preprint]
% is passed, so the real details live here and are revealed only at camera-ready.
\ifusetmlr
  \author{\name Youngmin Ko \email \\ \addr Pennsylvania State University}
\else
  \author{Anonymous authors\\Paper under double-blind review}
\fi

\date{}

\begin{document}
\maketitle

% TMLR requires an explicit first-page footnote disclosing LLM assistance.
\footnotetext[1]{Large language models were used as assistive tools in the
statistical re-analysis and in drafting this manuscript. All reported numbers
are produced by the released scripts and were verified against the committed
artifacts by the author, who takes full responsibility for the content.}

\begin{abstract}
Reports that point-level and segment-level metrics rank time-series anomaly
detectors differently have become a standard argument for evaluation reform.
These reports take the form of a \emph{rank-flip rate}: the fraction of model
pairs that two metrics order differently. We audit that statistic on a
\TsbadSeries-series, \TsbadModels-model recomputation of TSB-AD-M spanning
\TsbadCollections{} source collections, and on a \TabDatasets-dataset,
seven-detector grid. Our main finding is that evaluation series are treated as
independent when they are not, and that this invalidates the structural
explanations built on top of the statistic. The covariate most often invoked to
explain disagreement, the short-anomaly ratio $\alpha$, is constant within
\ConstAlpha{} of our \TsbadCollections{} collections: it is a collection label
rather than a series-level variable. Its rank correlation with the flip rate
falls from $\SerAlphaRho$ at the series level to $\ColAlphaRho$
($p = \ColAlphaP$) with collections as the unit, and to $\LocoAlphaMin$ when one
collection is removed. Interval estimates move the same way, from
$[\TsbadCiSeriesLo, \TsbadCiSeriesHi]$ resampling series to
$[\TsbadCiClusterLo, \TsbadCiClusterHi]$ clustering over collections; on the
six-dataset grid the clustered interval spans
$[\TabDeepAffCiClusterLo, \TabDeepAffCiClusterHi]$, which cannot support any
claim about magnitude. Two further omissions point the same way. Raw flip
percentages are reported without stating the chance level, which is $0.5$ and
not $0$: our observed \TsbadFlipAff{} against a permutation null of
\TsbadNull{} means the metrics agree on \TsbadAgreePct\% of pairs. And flips are
pooled across margins: among the \TsbadJointShareTwoZero\% of pairs that both
metrics separate by at least $0.20$, only \TsbadJointTwoZero{} are ordered
differently, so confident disagreement is roughly two percent rather than the
\TsbadFlipAff{} that gets reported. We apply
the audit to a composite metric we previously proposed ourselves and show it is
uninformative by construction at both ends of the regime it was built to span.
Mean segment duration and segment count do survive clustering. We give a
reporting protocol and release code that regenerates every number here from
committed artifacts.
\end{abstract}

\section{Introduction}

Time-series anomaly detection (TSAD) evaluation has been under sustained
criticism for a decade. Point-adjustment inflates scores to the point where
random detectors appear competitive \citep{kim2022towards}; widely used
benchmarks contain mislabeled, trivially separable, or unrepresentative anomalies
\citep{wu2023current, liu2024elephant}; and a series of replacement metrics has
been proposed to repair specific failures of point-wise scoring
\citep{tatbul2018precision, huet2022local, paparrizos2022vus, ghorbani2024pate}.

A particular kind of evidence recurs throughout this literature. One computes
two metrics over the same detector--dataset grid, counts the model pairs that
the two metrics order differently, and reports the fraction. We call this a
\emph{rank-flip rate}. It is rhetorically powerful: a statement that a third of
published rankings would reverse under a defensible alternative metric implies
that leaderboard positions carry little information. We used exactly this
argument in earlier work.

This paper asks what a rank-flip rate actually measures. Our finding is that
the statistic as reported is missing three ingredients, that each omission
pushes the interpretation in the same direction---toward overstating
instability---and that supplying them changes several conclusions in this
literature, including our own.

\paragraph{A stated chance level.} A flip rate is a fraction between $0$ and
$1$, and it is discussed as though $0$ were the reference point, so that any
positive value reads as evidence of disagreement. The correct reference is the
value the statistic takes when the two metrics share no information about model
quality, which is $0.5$. We are careful about the scope of this claim. The field
does baseline against chance elsewhere: \citet{yang2025problem} evaluate metric
\emph{values} under random detection, and \citet{lyu2026didwe} reports
inter-metric Kendall's $\tau$, which is itself chance-corrected. What is missing
is narrower and specific: raw flip \emph{percentages} are reported and
interpreted without the chance level being stated alongside them. The move of
supplying a permutation null for pairwise reversal rates has also been made
independently and concurrently outside our field---\citet{kendiukhov2026quantifying}
does exactly this for gene regulatory network benchmarking, obtaining a null of
$0.500$ against an observed $0.163$---which we take as corroboration rather than
as a contribution of ours.

\paragraph{A margin treatment.} Whether two metrics order a pair of models
differently is only meaningful if the models are actually distinguishable. Pairs
separated by a negligible margin flip at close to the chance rate by
construction, and they make up a substantial share of every grid. The margin
must moreover be taken on both metrics: conditioning only on the primary one
selects pairs it separates confidently while leaving the secondary margin free,
which is exactly where the flips turn out to be.

\paragraph{A cluster unit.} Evaluation series are routinely treated as
independent replicates. They are not: benchmark suites are assembled from source
collections, and series within a collection share a sensor layout, a labeling
convention, and an anomaly-generating process. This matters most for exactly the
covariates used to \emph{explain} disagreement, because those covariates tend to
be properties of the collection rather than of the series. When that is the case,
a structural explanation supported at $n = 180$ can be supported by, in effect, a
handful of independent observations.

\paragraph{Contributions.} Our primary contribution is the third item below;
the first two are corrections that are individually straightforward and, in the
case of the null, partly concurrent with work in another field.
\begin{enumerate}
\item We supply all three ingredients on two grids: a \TsbadSeries-series,
  \TsbadModels-model, \TsbadRows-row recomputation of TSB-AD-M spanning
  \TsbadCollections{} source collections, and a \TabDatasets-dataset,
  seven-detector study on TAB-style benchmarks (Section~\ref{sec:setup}).
\item We show that reported flip rates describe agreement once the chance level
  is stated, and that most of the residual is a margin artifact that survives
  only under one-sided conditioning
  (Sections~\ref{sec:null}--\ref{sec:margin}).
\item \textbf{We show that structure-conditioned explanations of disagreement
  are confounded with collection identity}, and that the standard thresholded
  short-anomaly ratio is not identifiable at the series level at all. We are not
  aware of a precedent for this analysis of rank-flip statistics, in TSAD or
  elsewhere (Section~\ref{sec:cluster}).
\item We give a worked failure case: a metric defined as an interpolation
  between two others is uninformative at both endpoints of the very regime it
  was introduced to span. The metric is one we proposed ourselves, and the
  analysis we published in support of it cannot come out any other way
  (Section~\ref{sec:interp}).
\item We identify two continuous covariates that do survive clustering, and
  give a reporting protocol (Sections~\ref{sec:survives}--\ref{sec:protocol}).
\end{enumerate}

We are not arguing that TSAD evaluation is sound. Two of the criticisms above
survive our audit intact, and a third---that benchmark anomalies are structurally
unlike the point-wise events that point metrics score---we confirm as a factual
matter. We are arguing that the specific statistic most often used to
\emph{quantify} the problem does not support the weight placed on it, and that
this is fixable.

\section{Setup}
\label{sec:setup}

\paragraph{Grids.} We use two evaluation grids. The first is a
\TabDatasets-dataset study on TAB-style benchmarks \citep{qiu2025tab}: SWaT
\citep{mathur2016swat, goh2017swat}, MSL and SMAP \citep{hundman2018detecting},
WADI \citep{ahmed2017wadi}, PSM, and SMD \citep{su2019robust}, evaluated with
five deep detectors---Anomaly Transformer \citep{xu2022anomaly}, TranAD
\citep{tuli2022tranad}, CATCH \citep{wu2025catch}, DADA \citep{shentu2025dada},
and MOMENT \citep{goswami2024moment}---plus two classical baselines, Isolation
Forest and LOF. The second is a replication on the TSB-AD-M evaluation list
\citep{liu2024elephant}, covering \TsbadSeries{} multivariate series and
\TsbadModels{} scored models for \TsbadRows{} model--series metric rows, drawn
from \TsbadCollections{} source collections.

\paragraph{Metrics.} We report \auc{} and segment-level affiliation F1
\citep{huet2022local} computed with \texttt{prts} at cardinality
\texttt{one} and flat bias, with thresholds set from the contamination rate of
the test partition following TAB's protocol \citep{qiu2025tab}. Contamination-rate
thresholding is itself an assumption with known sensitivity
\citep{masakuna2024impact, perini2023estimating}; our analysis conditions on it
throughout and we do not claim results are invariant to it.

\paragraph{Rank-flip rate.} For a group $g$ (a dataset, or a series) with model
scores $a_i$ under a primary metric and $b_i$ under a secondary metric, we count
unordered pairs $(i,j)$ with $a_i \neq a_j$ and $b_i \neq b_j$, and call a pair
\emph{flipped} when $\operatorname{sign}(a_i - a_j) \neq \operatorname{sign}(b_i - b_j)$.
The flip rate pools flipped pairs over comparable pairs. This matches the
definition used in the sources we are auditing, including our own.

\paragraph{Anomaly-duration taxonomy.} Anomaly segments are contiguous runs of
positive labels in the test partition, bucketed as Short ($\leq 5$ steps), Medium
($6$--$50$), or Long ($> 50$). The structure weight in common use is
$\alpha = 1 - N_{\text{short}} / N_{\text{total}}$. On the
\TabDatasets-dataset grid, only PSM ($\alpha = \AlphaPSM$) and SMD
($\alpha = \AlphaSMD$) have $\alpha < 1$; SWaT, MSL, SMAP, and WADI have no Short
segments at all. That last fact is a property of the labels and is not disputed
here.

\paragraph{Reproducibility.} Every number in this paper is emitted as a \LaTeX{}
macro by a script that reads the committed JSON artifacts, so the manuscript
cannot drift from the analysis. Rerunning three scripts regenerates all
artifacts and all numbers.

\section{Against a null, the metrics mostly agree}
\label{sec:null}

We build the null by permuting the secondary metric across models within each
group, which destroys any association with the primary metric while preserving
both marginal distributions and all group sizes. Under independent rankings the
expected flip rate is $0.5$.

Table~\ref{tab:null} gives the result. On TSB-AD-M the observed
\auc{}-versus-\afff{} rate is \TsbadFlipAff{} against a null of \TsbadNull{}
(sd \TsbadNullSd{}, \TsbadNullPerm{} permutations). On the
\TabDatasets-dataset grid the deep-model rate is
\TabDeepAffFlips/\TabDeepAffPairs{} $=$ \TabDeepAffRate{} against
\TabDeepAffNull{}, and the seven-model rate is
\TabAllAffFlips/\TabAllAffPairs{} $=$ \TabAllAffRate{} against
\TabAllAffNull{}. Every observed rate is far below its null.

\begin{table}[t]
\centering
\caption{Rank-flip rates against a random-ranking null. Reported rates are
routinely discussed against an implicit reference of zero; the informative
reference is the null column. All rates are \auc{} versus \afff{}.}
\label{tab:null}
\begin{tabular}{llrrrr}
\toprule
Grid & Models & Flips / pairs & Rate & Null & Agreement \\
\midrule
TSB-AD-M & \TsbadModels{} & --- & \TsbadFlipAff{} & \TsbadNull{} & \TsbadAgreePct\% \\
TAB & 5 deep & \TabDeepAffFlips/\TabDeepAffPairs & \TabDeepAffRate{} & \TabDeepAffNull{} & \TabDeepAffAgreePct\% \\
TAB & 7 (with classical) & \TabAllAffFlips/\TabAllAffPairs & \TabAllAffRate{} & \TabAllAffNull{} & \TabAllAffAgreePct\% \\
\bottomrule
\end{tabular}
\end{table}

The consequence is a reframing rather than a refutation. ``\TabAllAffAgreePct\%
of pairwise rankings are preserved across a point metric and a segment metric''
and ``\TabAllAffRate{} of pairwise rankings flip'' are the same measurement. The
first is the honest summary; the second, which is how such numbers are presented,
invites the reader to supply a reference point of zero.

We claim no novelty for this step. \citet{kendiukhov2026quantifying} applies the
same permutation construction to ranking reversals in gene regulatory network
benchmarking and reaches the analogous conclusion, that a headline reversal rate
of $0.163$ against a null of $0.500$ indicates substantial shared ranking
structure. Within TSAD, chance-corrected rank agreement is also already
available in the form of Kendall's $\tau$ \citep{lyu2026didwe}. The point is
that the raw percentage is what circulates, and it circulates without its
baseline.

\section{Most of the disagreement is a margin artifact}
\label{sec:margin}

Two models separated by $0.001$ \auc{} are not meaningfully ordered by \auc{},
so a metric that orders them differently has not reversed a finding. Because
such pairs are numerous, they can carry a flip rate on their own.

Table~\ref{tab:margin} stratifies the TSB-AD-M pairs by $|\Delta\auc{}|$. The
pattern is monotone and steep: pairs separated by less than $0.01$ flip at
\TsbadGapZeroZeroToZeroOne{}, essentially the chance rate, while pairs separated
by at least $0.20$ flip at \TsbadGapTwoZeroToInf{}. Excluding every pair within
$0.05$ \auc{} retains \TsbadDropShareZeroFive\% of pairs and gives
\TsbadDropZeroFive{}.

\paragraph{One-sided conditioning is not enough.} It is tempting to stop here
and report that roughly one pair in five still reverses at a margin no one would
call ambiguous. That reading is wrong, and the error is instructive. Conditioning
on $|\Delta\auc{}|$ selects pairs that the \emph{primary} metric separates
confidently while saying nothing about whether \afff{} separates them at all, so
the surviving flips concentrate where the secondary margin is negligible. The
symmetric stratification confirms the asymmetry: pairs with
$|\Delta\afff{}| < 0.02$ flip at \TsbadSecZeroZeroToZeroTwo{}---the chance
rate---and pairs with $|\Delta\afff{}| \geq 0.40$ flip at
\TsbadSecFourZeroToInf{}.

Requiring both metrics to separate the pair gives the two-sided statistic
(Table~\ref{tab:joint}). Among the \TsbadJointShareTwoZero\% of pairs that
\emph{both} metrics separate by at least $0.20$, only \TsbadJointTwoZero{} are
ordered differently. The corresponding figures at joint margins of $0.05$ and
$0.10$ are \TsbadJointZeroFive{} and \TsbadJointOneZero{}.

We regard the two-sided number as the answer to the question the literature is
actually asking. ``How often do two metrics confidently disagree about which of
two detectors is better?'' is not \TsbadFlipAff{}, and it is not
\TsbadGapTwoZeroToInf{}. On this grid it is about two percent.

\begin{table}[t]
\centering
\caption{TSB-AD-M flip rate by \auc{} margin. Disagreement falls steeply with
margin, but conditioning on the primary metric alone is one-sided; see
Table~\ref{tab:joint}.}
\label{tab:margin}
\begin{tabular}{lrrr}
\toprule
$|\Delta\auc{}|$ & Pairs & Share & Flip rate \\
\midrule
$< 0.01$        & \TsbadGapPairsZeroZeroToZeroOne   & \TsbadGapShareZeroZeroToZeroOne\%   & \TsbadGapZeroZeroToZeroOne \\
$0.01$--$0.02$  & \TsbadGapPairsZeroOneToZeroTwo    & \TsbadGapShareZeroOneToZeroTwo\%    & \TsbadGapZeroOneToZeroTwo \\
$0.02$--$0.05$  & \TsbadGapPairsZeroTwoToZeroFive   & \TsbadGapShareZeroTwoToZeroFive\%   & \TsbadGapZeroTwoToZeroFive \\
$0.05$--$0.10$  & \TsbadGapPairsZeroFiveToOneZero   & \TsbadGapShareZeroFiveToOneZero\%   & \TsbadGapZeroFiveToOneZero \\
$0.10$--$0.20$  & \TsbadGapPairsOneZeroToTwoZero    & \TsbadGapShareOneZeroToTwoZero\%    & \TsbadGapOneZeroToTwoZero \\
$\geq 0.20$     & \TsbadGapPairsTwoZeroToInf        & \TsbadGapShareTwoZeroToInf\%        & \TsbadGapTwoZeroToInf \\
\bottomrule
\end{tabular}
\end{table}

\begin{table}[t]
\centering
\caption{Two-sided margin conditioning on TSB-AD-M. Left: stratified by the
secondary metric's margin. Right: both margins required to exceed a threshold.
Confident disagreement is rare.}
\label{tab:joint}
\begin{tabular}{lrr@{\hskip 2.5em}lrr}
\toprule
$|\Delta\afff{}|$ & Pairs & Flip & Both margins $\geq$ & Pairs & Flip \\
\midrule
$< 0.02$       & \TsbadSecPairsZeroZeroToZeroTwo  & \TsbadSecZeroZeroToZeroTwo  & $0.05$ & \TsbadJointPairsZeroFive & \TsbadJointZeroFive \\
$0.02$--$0.05$ & \TsbadSecPairsZeroTwoToZeroFive  & \TsbadSecZeroTwoToZeroFive  & $0.10$ & \TsbadJointPairsOneZero  & \TsbadJointOneZero \\
$0.05$--$0.10$ & \TsbadSecPairsZeroFiveToOneZero  & \TsbadSecZeroFiveToOneZero  & $0.20$ & \TsbadJointPairsTwoZero  & \TsbadJointTwoZero \\
$0.10$--$0.20$ & \TsbadSecPairsOneZeroToTwoZero   & \TsbadSecOneZeroToTwoZero   &        &                          & \\
$0.20$--$0.40$ & \TsbadSecPairsTwoZeroToFourZero  & \TsbadSecTwoZeroToFourZero  &        &                          & \\
$\geq 0.40$    & \TsbadSecPairsFourZeroToInf      & \TsbadSecFourZeroToInf      &        &                          & \\
\bottomrule
\end{tabular}
\end{table}

\begin{table}[t]
\centering
\caption{TSB-AD-M flip rate by \auc{} margin. Disagreement falls steeply with
margin but does not vanish.}
\label{tab:margin}
\begin{tabular}{lrrr}
\toprule
$|\Delta\auc{}|$ & Pairs & Share & Flip rate \\
\midrule
$< 0.01$        & \TsbadGapPairsZeroZeroToZeroOne   & \TsbadGapShareZeroZeroToZeroOne\%   & \TsbadGapZeroZeroToZeroOne \\
$0.01$--$0.02$  & \TsbadGapPairsZeroOneToZeroTwo    & \TsbadGapShareZeroOneToZeroTwo\%    & \TsbadGapZeroOneToZeroTwo \\
$0.02$--$0.05$  & \TsbadGapPairsZeroTwoToZeroFive   & \TsbadGapShareZeroTwoToZeroFive\%   & \TsbadGapZeroTwoToZeroFive \\
$0.05$--$0.10$  & \TsbadGapPairsZeroFiveToOneZero   & \TsbadGapShareZeroFiveToOneZero\%   & \TsbadGapZeroFiveToOneZero \\
$0.10$--$0.20$  & \TsbadGapPairsOneZeroToTwoZero    & \TsbadGapShareOneZeroToTwoZero\%    & \TsbadGapOneZeroToTwoZero \\
$\geq 0.20$     & \TsbadGapPairsTwoZeroToInf        & \TsbadGapShareTwoZeroToInf\%        & \TsbadGapTwoZeroToInf \\
\bottomrule
\end{tabular}
\end{table}

A related concern is that flips are produced by detectors that do not
discriminate at all. On the \TabDatasets-dataset grid we drop every
model--dataset row with $|\auc{} - 0.5| < \delta$ and recompute. The rate is
\NrRateZeroZero{} at $\delta = 0$, then \NrRateZeroTwo{}, \NrRateZeroFive{}, and
\NrFlipsOneZero/\NrPairsOneZero{} $=$ \NrRateOneZero{} at
$\delta = 0.02, 0.05, 0.10$. Removing uninformative detectors does not remove
the disagreement.

Two cautions about this check. The surviving sample is small---%
\NrPairsOneZero{} pairs at $\delta = 0.10$---and the clustered interval is
correspondingly wide, $[\NrCiLoOneZero, \NrCiHiOneZero]$. And excluding by
measured discriminative power is not the same as excluding by model family:
Isolation Forest is the strongest detector by \auc{} on four of the six
datasets, while Anomaly Transformer sits within $0.03$ of chance on all six, so
the $\delta = 0.10$ exclusion drops \NrDeepOneZero{} deep rows against
\NrClsOneZero{} classical ones. Pairs that mix a deep and a classical detector do
flip more often (\ClsDCFlips/\ClsDCPairs{} $=$ \ClsDCRate{}, versus
\ClsDDFlips/\ClsDDPairs{} $=$ \ClsDDRate{} for deep--deep pairs), but this traces
to LOF specifically rather than to classical methods as a class.

\section{Evaluation series are not independent}
\label{sec:cluster}

TSB-AD-M draws its \TsbadSeries{} series from \TsbadCollections{} source
collections. Pooled flip rates differ sharply across them, from \CollMinRate{}
on \CollMinName{} to \CollMaxRate{} on \CollMaxName{}. Treating series as
independent replicates therefore overstates precision, and---more
seriously---creates a confound for any covariate that is a property of the
collection.

The short-anomaly ratio is such a covariate. Table~\ref{tab:ident} shows that
$\alpha$ is constant within \ConstAlpha{} of \TsbadCollections{} collections, and
within \ConstBigAlpha{} of the \NBigCollections{} collections that contribute at
least eight series. It is, operationally, a collection label. Mean segment
duration and segment count are not: both vary within every one of the large
collections.

\begin{table}[t]
\centering
\caption{Within-collection variation of structural covariates on TSB-AD-M.
A covariate that is constant within collections cannot be identified at the
series level.}
\label{tab:ident}
\begin{tabular}{lrr}
\toprule
Covariate & Constant within & Constant within large \\
          & (of \TsbadCollections{} collections) & (of \NBigCollections{} with $\geq 8$ series) \\
\midrule
$\alpha = 1 - N_{\text{short}}/N_{\text{total}}$ & \ConstAlpha{} & \ConstBigAlpha{} \\
Mean segment duration & \ConstDur{} & \ConstBigDur{} \\
Segment count & \ConstNseg{} & \ConstBigNseg{} \\
\bottomrule
\end{tabular}
\end{table}

The consequences are in Table~\ref{tab:pred}. At the series level $\alpha$
correlates with the flip rate at $\rho = \SerAlphaRho{}$ with a permutation
$p$ at the resolution floor of the test---an apparently decisive result. At the
collection level the same association is $\ColAlphaRho{}$ with
$p = \ColAlphaP{}$. Leave-one-collection-out is equally stark: removing
\LocoAlphaWorst{} alone takes the series-level correlation to
\LocoAlphaMin{}. The $\alpha = 0$ stratum on this grid consists of eleven
\LocoAlphaWorst{} series and one other, so the low-$\alpha$ regime is one
collection wearing the clothes of twelve observations.

The same holds on the \TabDatasets-dataset grid, independently. Only
\TabAlphaLtOne{} of \TabDatasets{} datasets have $\alpha < 1$, so the entire
structural claim rests on them; and of the \TabDeepAffFlips{} deep-model flips,
\DeepPSMFlips{} come from PSM and \DeepSWaTFlips{} from SWaT. SWaT has
$\alpha = \AlphaSWaT{}$---the regime in which the structural account predicts
the \emph{fewest} flips.

Finally, spread within the $\alpha = 1$ regime exceeds spread across $\alpha$
regimes: among collections with $\alpha = 1$ throughout, pooled flip rates run
from \AOneCollMinRate{} (\AOneCollMinName{}) to \AOneCollMaxRate{}
(\AOneCollMaxName{}), a wider range than the $\alpha$-stratified contrast
itself.

\begin{table}[t]
\centering
\caption{Spearman correlation between structural covariates and the flip rate,
at series level and with collections as the unit, plus the range under
leave-one-collection-out. Series-level $p$-values are permutation-based and
report the resolution floor of the test. $\alpha$ is the only covariate that
does not survive.}
\label{tab:pred}
\begin{tabular}{lrrrrl}
\toprule
Covariate & Series $\rho$ & Coll.\ $\rho$ & Coll.\ $p$ & LOCO range & Worst drop \\
\midrule
$\alpha$              & \SerAlphaRho & \ColAlphaRho & \ColAlphaP & $[\LocoAlphaMin, \LocoAlphaMax]$ & \LocoAlphaWorst \\
Mean segment duration & \SerDurRho   & \ColDurRho   & \ColDurP   & $[\LocoDurMin, \LocoDurMax]$     & \LocoDurWorst \\
Segment count         & \SerNsegRho  & \ColNsegRho  & \ColNsegP  & $[\LocoNsegMin, \LocoNsegMax]$   & \LocoNsegWorst \\
Anomaly density       & \SerDensRho  & \ColDensRho  & \ColDensP  & $[\LocoDensMin, \LocoDensMax]$   & \LocoDensWorst \\
Series length         & \SerLenRho   & \ColLenRho   & \ColLenP   & $[\LocoLenMin, \LocoLenMax]$     & \LocoLenWorst \\
\bottomrule
\end{tabular}
\end{table}

\paragraph{Interval estimates.} The same non-independence affects uncertainty.
Resampling series gives a 95\% interval of
$[\TsbadCiSeriesLo, \TsbadCiSeriesHi]$ for the TSB-AD-M flip rate; a cluster
bootstrap over collections gives $[\TsbadCiClusterLo, \TsbadCiClusterHi]$, about
three times wider. On the \TabDatasets-dataset grid the effect is severe enough
to be disqualifying: the deep-model interval is
$[\TabDeepAffCiPairLo, \TabDeepAffCiPairHi]$ at pair level but
$[\TabDeepAffCiClusterLo, \TabDeepAffCiClusterHi]$ clustering over the six
datasets. A study of that size cannot resolve a flip rate to better than a
range spanning most of the unit interval, and should not be used to argue about
its magnitude.

\section{A worked failure case: interpolation metrics at their endpoints}
\label{sec:interp}

One response to metric disagreement is to combine metrics. In earlier work we
proposed such a composite, which we called \sae{}, with the structure weight
$\alpha$ selecting the blend:
\begin{equation}
\sae{}(D, M) = (1 - \alpha(D))\,\auc{}(M) + \alpha(D)\,\afff{}(M),
\qquad \alpha(D) = 1 - \frac{N_{\text{short}}(D)}{N_{\text{total}}(D)} .
\label{eq:sae}
\end{equation}
We then stratified the \auc{}-versus-\sae{} flip rate by $\alpha$ and reported
the resulting gradient---\StratSaeAZero{}, \StratSaeAMid{}, \StratSaeAOne{}
across $\alpha = 0$, $0 < \alpha < 1$, and $\alpha = 1$---as confirming a
structural regime split.

It confirms nothing. Equation~\ref{eq:sae} makes \sae{} identically \auc{} at
$\alpha = 0$ and identically \afff{} at $\alpha = 1$; on our artifact the
maximum absolute differences are \IdentAlphaZero{} and \IdentAlphaOne{}
respectively. The $\alpha = 0$ stratum therefore compares \auc{} with itself and
contributes \StratSaeAZeroFlips{} flips out of \StratSaeAZeroPairs{} by
definition, and the $\alpha = 1$ stratum is bit-identical to the
\auc{}-versus-\afff{} comparison (\StratSaeAOneFlips/\StratSaeAOnePairs{} in
both). The gradient is linear interpolation between a self-comparison and a
fixed comparison. No configuration of the data could have produced anything
else.

This generalizes past our own metric. Any composite defined as a convex
combination of two metrics, with a weight that reaches its endpoints on real
data, will exhibit a monotone ``regime'' pattern when stratified by that weight,
whether or not the underlying structural hypothesis is true. Such a stratification
is a description of the metric's definition, not a finding about benchmarks. The
diagnostic is cheap: check whether the extreme strata of the weight reduce the
composite to one of its constituents, and if they do, report the constituent
comparison instead.

Doing that here gives strata of \StratAffAZero{}, \StratAffAMid{}, and
\StratAffAOne{}. These are metric-independent and still monotone---but as
Section~\ref{sec:cluster} shows, the gradient is carried by collection identity,
and $\alpha$ does not survive as a predictor.

\section{What survives}
\label{sec:survives}

The structural hypothesis is not dead; its standard operationalization is.
Mean segment duration correlates with the flip rate at $\SerDurRho{}$ at the
series level and $\ColDurRho{}$ at the collection level, and stays within
$[\LocoDurMin, \LocoDurMax]$ under leave-one-collection-out. Segment count
behaves symmetrically at $\SerNsegRho{}$ and $\ColNsegRho{}$. Neither collapses
the way $\alpha$ does, and neither is constant within the large collections.

We state the limits plainly. With \TsbadCollections{} collections the
collection-level tests are underpowered: the permutation $p$-values are
$\ColDurP{}$ and $\ColNsegP{}$, and we do not claim significance for them. What
we claim is the comparison. Under identical treatment, $\alpha$ moves from
$\SerAlphaRho{}$ to $\ColAlphaRho{}$ and collapses under leave-one-out, whereas
duration and segment count retain their sign and roughly their magnitude at both
levels. That contrast is the evidence; the individual collection-level
coefficients are not.

Note also that these covariates are not independent of one another. Mean segment
duration is algebraically determined by anomaly density, series length and
segment count, so a regression that conditions on all three cannot say anything
further about duration. Reporting several one-at-a-time associations, as in
Table~\ref{tab:pred}, is the honest presentation at this sample size.

\section{A reporting protocol}
\label{sec:protocol}

The following costs no additional experiments and would have caught every issue
above.

\begin{enumerate}
\item \textbf{Report the null.} State the random-ranking baseline next to any
  metric-disagreement statistic. For pairwise flip rates it is $0.5$; obtain it
  by permuting one metric across models within each evaluation group.
\item \textbf{Stratify by margin.} Report the flip rate as a function of the
  primary metric's gap, or at minimum report it with near-ties excluded and say
  what fraction of pairs that removes.
\item \textbf{Cluster at the source.} Use the source collection or dataset, not
  the series, as the resampling unit, and report the clustered interval. Where
  the two intervals differ materially, report both.
\item \textbf{Check identifiability before conditioning.} Before explaining
  disagreement with a covariate, check whether that covariate varies within
  clusters. If it does not, a series-level association is a between-collection
  association with an inflated $n$.
\item \textbf{Leave one collection out.} Report the range of any structural
  association under removal of each collection in turn. A result carried by one
  collection should be presented as such.
\item \textbf{Do not stratify a composite by its own weight.} If a metric
  reduces to a constituent at the extremes of a weight, stratifying by that
  weight measures the definition.
\end{enumerate}

\section{Related work}

Our object of study is a statistic, not a metric, so the closest work is
methodological rather than metric-proposing.

\paragraph{Critiques of TSAD evaluation.} \citet{kim2022towards} showed that
point-adjustment lets random scores outperform published detectors, which is the
canonical demonstration that a TSAD metric can be uninformative.
\citet{wu2023current} and \citet{liu2024elephant} argue that benchmark datasets
themselves are flawed, mislabeled, or trivially solvable, and
\citet{schmidl2022anomaly} and \citet{boniol2024dive} survey the resulting
landscape. \citet{lyu2026didwe} stress-tests post-point-adjustment replacement
metrics and finds that several, including affiliation-based and ROC-based
scores, inflate steeply under best-of-$N$ seed selection while PR-based metrics
stay flat. That result is complementary to ours and sharpens it: the metric on
one side of our comparison has a known instability under selection, which is a
further reason not to read a raw flip count as a property of the benchmark.

\paragraph{Replacement metrics.} Range-based precision and recall
\citep{tatbul2018precision}, affiliation metrics \citep{huet2022local}, VUS
\citep{paparrizos2022vus}, PATE \citep{ghorbani2024pate}, DQE
\citep{li2026dqe}, and confidence-consistency evaluation \citep{zhong2025cce}
each repair a specific deficiency of point-wise scoring.
\citet{wagner2025formally} approaches the space axiomatically, specifying
properties a metric ought to satisfy, and \citet{yang2025problem} organize the
resulting metric space by the evaluation problem each metric targets---including
a dimension for robustness against random-score inflation, which is chance
baselining applied to metric values. Our contribution is orthogonal to all of
these: whichever metric one adopts, comparisons \emph{between} metrics are
currently reported without the ingredients above.

\paragraph{Ranking instability as a statistic.} Quantifying how far rankings
move under a change of protocol is now common practice across ML evaluation,
from preference-data perturbations in language model leaderboards
\citep{dropping2025preferences} to protocol axes in computational biology
\citep{kendiukhov2026quantifying}. The latter is the closest methodological
neighbour to our Section~\ref{sec:null}, and reaches the same conclusion in a
different domain. It does not address clustering; its own limitations section
lists confounding with tissue identity as unresolved, which is the analogue of
the problem we take up in Section~\ref{sec:cluster}.

\paragraph{Statistical practice in benchmark evaluation.} The concern that
evaluation units are not independent replicates is old outside machine learning
\citep{hurlbert1984pseudoreplication}. Within it, the argument that point
estimates on small benchmark suites mislead without interval estimates has been
made forcefully in reinforcement learning \citep{agarwal2021precipice} and for
language model evaluations \citep{miller2024errorbars}, and
\citet{rochner2025rethink} argue on general grounds that anomaly detection
benchmarks should be organized into scenario groups sharing relevant
characteristics. Ours is the quantitative counterpart to that last argument:
grouping is not only a taxonomic convenience but the difference between a
structural association of $\SerAlphaRho{}$ and one of $\ColAlphaRho{}$. Relative
to \citet{agarwal2021precipice} and \citet{miller2024errorbars}, which resample
runs or items as exchangeable units, our point is that the exchangeable unit
here is the source collection, and that identifying it changes which covariates
remain admissible at all.

\section{Limitations}

Our conclusions are conditioned on a specific metric pair, \auc{} versus \afff{}
at cardinality \texttt{one} with contamination-rate thresholding. Other pairs
may behave differently, and the thresholding rule is known to matter
\citep{masakuna2024impact}. The TSB-AD-M slice follows file-ID order and is
model-coverage-limited rather than a complete benchmark run, so it should be read
as a replication of the statistical phenomenon and not as a TSB-AD leaderboard
claim. With \TsbadCollections{} collections and \TabDatasets{} datasets, the
cluster-level analyses that carry our main argument are themselves
underpowered---which is part of the point, but does mean our positive claims
about duration and segment count are directional rather than established. We
have not shown that any published conclusion in the TSAD literature is wrong as
a result of these omissions; we have shown that the evidence offered for one
class of conclusion, including ours, does not have the strength attributed to it.

\section{Conclusion}

A rank-flip rate is a compact and persuasive statistic, and it is currently
reported in a way that cannot be interpreted. Supplied with a null, the flip
rates in this literature and in our own prior work describe metrics that
substantially agree. Stratified by margin, a real residual disagreement remains
that near-ties do not explain. Clustered at the source collection, the structural
account of that disagreement does not survive, and neither does the composite
metric we built on it. What we can defend is narrower than what we previously
claimed, and we think it is more useful: point and segment metrics confidently
disagree on about two percent of model pairs, what disagreement there is varies
across benchmark collections to a degree that dwarfs any anomaly-structure
effect we can identify, and the structural covariates that do track it are
continuous ones the field has not been using.

\subsubsection*{Broader Impact Statement}

This work audits an evaluation statistic and proposes reporting practice; it
introduces no model, no dataset, and no deployable artifact. The direct effect
we anticipate is that some published claims about metric instability in
time-series anomaly detection will be restated more conservatively, including
our own. We see one indirect risk worth naming. Showing that a widely cited
instability figure overstates disagreement could be read as licence to keep
reporting single point-wise metrics on benchmarks whose anomalies are sustained
segments. That reading would be wrong, though our results make the point less forcefully
than the ones they replace: the labeling mismatch documented in prior work is
untouched by anything we show, and a metric can be misaligned with a task even
when it happens to rank detectors much as another metric does. The
benchmarks we use are public industrial and infrastructure datasets containing
no personal data.

\subsubsection*{Reproducibility statement}

All results are recomputed from committed JSON artifacts by three scripts, and
every numeric claim in this manuscript is generated into \texttt{numbers.tex} by
a fourth, so the text cannot diverge from the analysis. The statistical helpers
are implemented without a scipy dependency and are unit-tested against
hand-computed values. The repository includes the derived metric rows, the
taxonomy summary, and the reproduction commands; it does not redistribute the
source benchmark datasets, for which access instructions are provided.

\ifusetmlr
  \bibliographystyle{tmlr}
\else
  \bibliographystyle{plainnat}
\fi
\bibliography{refs}

\end{document}
